\\section{AI-Assisted Reanalysis Pipeline}
\label{AppendixC}

This appendix documents the automated pipeline used to reanalyze 53 TWFE studies with modern DiD estimators. The pipeline is a nine-agent system in which a language-model orchestrator (Claude Sonnet~4, Anthropic~2025, accessed through Claude Code) coordinates deterministic computational agents through structured file-based communication. The architecture follows the framework proposed by \textcite{xu2026scaling}, who evaluate a similar design on 92 instrumental-variable studies (215 total specifications across seven pipeline stages) and report 87\% end-to-end success. The present project adapts that framework to the heterogeneity of DiD replication packages.

The remainder of this appendix is organized as follows. Section~\ref{sec:appc_design} states the design principles. Section~\ref{sec:appc_architecture} introduces the nine-agent architecture. Sections~\ref{sec:appc_execution}--\ref{sec:appc_custodian} describe each agent family in turn. Section~\ref{sec:appc_metadata} specifies the metadata contract that connects the orchestrator to the execution layer. Section~\ref{sec:appc_calibration} describes the validation protocol. Section~\ref{sec:appc_xuyang} compares the architecture to \textcite{xu2026scaling}. Section~\ref{sec:appc_directory} documents the directory layout, and Section~\ref{sec:appc_replication} gives replication instructions.


\\subsection{Design principles}
\label{sec:appc_design}

The pipeline rests on four principles. First, reasoning and execution are separated. The language-model orchestrator reads papers and makes specification decisions, and a fixed R script handles all estimation. The two layers communicate exclusively through a structured JSON file (\texttt{metadata.json}), so no analytical choice is hidden inside code modifications.

Second, the template is fixed. The R template \texttt{did\_analysis\_template.R} is never modified per article. Given identical metadata and data, the template produces identical results, which eliminates a large class of replication failures tied to ad-hoc script edits. All article-specific logic lives in \texttt{metadata.json}.

Third, institutional memory is version-controlled. Implementation failures are recorded in a knowledge base (\texttt{knowledge/failure\_patterns.md}). The orchestrator consults this file before profiling each new article, so the system does not relearn resolved pitfalls.

Fourth, inter-agent communication is file-based. An interrupted run can be resumed from the last completed stage without loss of state, and every artefact produced by any agent is reproducible from the inputs visible on disk.


\\subsection{Architecture}
\label{sec:appc_architecture}

The system is organized as a single orchestrator coordinating nine agents arranged in four layers (Figure~\ref{fig:pipeline_arch}). Three execution agents (Profiler, Analyst, Reporter) produce the numerical results. Six audit agents verify them. Five are methodology reviewers (TWFE, CS-DID, Bacon, HonestDiD, de~Chaisemartin--d'Haultf\oe{}uille). One is a fidelity reviewer (paper-auditor). One is an orchestrator (skeptic) that synthesises their verdicts into a credibility rating. Two custodial agents (pattern-curator and repo-custodian) maintain the knowledge base and the git repository.

\begin{figure}[H]
\centering
\begin{tcolorbox}[colback=gray!5, colframe=gray!60, width=\textwidth, title=Nine-agent architecture]
\footnotesize
\begin{verbatim}
  Orchestrator (Claude Sonnet 4, governed by CLAUDE.md)
     |
     |-- reads ---> Knowledge base (failure_patterns.md)
     |              Skill contracts  (SKILL_profiler.md, _analyst.md, _reporter.md)
     |
     |-- EXECUTION LAYER (3 agents) ---------------------------------------
     |   Profiler   : PDF + code + data  ->  metadata.json
     |   Analyst    : metadata.json + data -> results.csv + event_study.pdf
     |   Reporter   : results + metadata   -> report.md
     |
     |-- AUDIT LAYER (5 methodology reviewers, read-only) ------------------
     |   twfe-reviewer           : FE structure, clustering, Pattern-50 check
     |   csdid-reviewer          : att_gt arguments, xformla, base period
     |   bacon-reviewer          : applicability + TvT weight interpretation
     |   honestdid-reviewer      : event-study inputs + M-bar breakdown values
     |   dechaisemartin-reviewer : advises on reversible / continuous treatment
     |
     |-- FIDELITY LAYER (read-only) ---------------------------------------
     |   paper-auditor : compares stored coefficient to published number
     |                   (EXACT / WITHIN_TOL / WARN / FAIL / NOT_APPLICABLE)
     |
     |-- ORCHESTRATION LAYER ---------------------------------------------
     |   skeptic       : runs the 6 audits in parallel and emits
     |                   HIGH / MODERATE / LOW + design-credibility finding
     |
     '-- CUSTODIAL LAYER -------------------------------------------------
         pattern-curator : maintains failure_patterns.md
         repo-custodian  : owns git, full_sync.sh, and the 24-check audit
\end{verbatim}
\end{tcolorbox}
\caption{The nine-agent architecture}
\label{fig:pipeline_arch}
\vspace{4pt}
{\footnotesize\textit{Notes:} The orchestrator coordinates three execution agents, five methodology reviewers, one fidelity reviewer, one orchestrator sub-agent (skeptic), and two custodial agents.\par}
\end{figure}


\\subsection{Execution agents}
\label{sec:appc_execution}

The three execution agents produce the numerical results that the downstream audit layer evaluates.

\\subsubsection*{Profiler}

The Profiler reads the article PDF (\texttt{pdf/[id].pdf}) and the replication package (\texttt{data/[id]/}) and writes \texttt{results/[id]/metadata.json}. It identifies the outcome variable, treatment indicator, unit of observation, time period, and whether treatment adoption is staggered or single-date. It then reads the replication code (\texttt{.do} or \texttt{.R} files) to extract variable names, sample restrictions, control variables, and clustering choices, and inspects the data file to confirm that variable names match (case-sensitive in R).

When multiple specifications are reported, the Profiler applies a hierarchical selection protocol. A specification explicitly labeled ``main,'' ``preferred,'' or ``baseline'' is selected first. If no such label exists, the specification with the largest $N$ is selected. If still tied, the specification with the fewest time-varying controls is selected. If still tied, the first specification presented is selected. The Profiler then fills all fields of \texttt{metadata.json} (Section~\ref{sec:appc_metadata}) and validates internal consistency. In particular, \texttt{run\_csdid\_nyt} and \texttt{run\_bacon} are set to \texttt{true} if and only if treatment timing is staggered, and \texttt{has\_event\_study} reflects the original paper.

\\subsubsection*{Analyst}

The Analyst takes \texttt{metadata.json} and the raw data as input and produces \texttt{results.csv} and, where applicable, \texttt{event\_study.pdf}. The entire estimation runs through a single command.
\begin{quote}
\texttt{Rscript code/analysis/did\_analysis\_template.R [id]}
\end{quote}
The template performs estimations conditional on flags set by the Profiler. TWFE is always estimated via \texttt{fixest::feols} with unit and time fixed effects and clustered standard errors. CS-DID with never-treated controls is always estimated via \texttt{did::att\_gt} with \texttt{control\_group = "nevertreated"}, aggregated to a scalar ATT via \texttt{aggte}. When treatment adoption is staggered, the template additionally runs CS-DID with not-yet-treated controls and a Bacon decomposition (\texttt{bacondecomp::bacon}). If the original paper reports an event study, Sun--Abraham estimates are obtained via \texttt{fixest::feols} with \texttt{sunab()}, and Gardner and BJS estimates via \texttt{did2s::did2s}.

\\subsubsection*{Reporter}

The Reporter reads \texttt{results.csv}, \texttt{event\_study.pdf}, and \texttt{metadata.json} to produce a structured report (\texttt{results/[id]/report.md}) covering six sections. The six are article summary, specification, point estimates, comparison with the original finding, event-study assessment, and conclusion. The conclusion classifies each reanalysis as \emph{confirmed} (CS-DID and TWFE agree in sign and significance, with magnitude difference below 20\%), \emph{partial} (same sign and significance but magnitude differs by 20\% or more, or a marginal significance change), or \emph{challenged} (sign reversal, or one estimator significant while the other is not).


\\subsection{Methodology audit agents}
\label{sec:appc_reviewers}

Five read-only methodology reviewers evaluate the Analyst's output against best-practice implementation checks. Each owns one estimator family and emits a PASS / WARN / FAIL verdict plus a structured report.

\\subsubsection*{\texttt{twfe-reviewer}}

Validates the FE structure, clustering, weights, and covariates of the TWFE regression against the paper's own specification. It applies in particular the Pattern-50 sample-mismatch invariant. The same sample must enter the TWFE and CS-DID branches unless an explicit \texttt{cs\_sample\_filter} or \texttt{cs\_construct\_vars} is declared. Silent sample mismatches of the Kresch-(2020) class are caught at this stage.

\\subsubsection*{\texttt{csdid-reviewer}}

Audits \texttt{did::att\_gt} usage on four margins, namely control-group choice, \texttt{xformla}, base period, and balancing flags. It reports whether CS-DID was applied under the preconditions required by \textcite{callaway2021difference}.

\\subsubsection*{\texttt{bacon-reviewer}}

Determines whether \texttt{bacondecomp::bacon} is formally applicable (balanced panel with staggered adoption). When applicable, it interprets the treated-vs-treated (TvT) weight share as an attenuation-risk signal consistent with \textcite{goodman2021difference}.

\\subsubsection*{\texttt{honestdid-reviewer}}

Checks the event-study inputs fed into \texttt{HonestDiD}, computes the $\bar M$ breakdown values under Rambachan--Roth relative-magnitude restrictions, and translates them into the D-ROBUST, D-MODERATE, D-FRAGILE, and D-BROKEN scale documented in Appendix~\ref{AppendixB}.

\\subsubsection*{\texttt{dechaisemartin-reviewer}}

Advises when the \texttt{DIDmultiplegtDYN} estimator would add information. This applies when treatment is reversible or continuous, or when TWFE and CS-DID disagree strongly and treatment reversal is suspected.


\\subsection{Fidelity audit agent}
\label{sec:appc_fidelity}

The \texttt{paper-auditor} agent answers a separate, narrower question. Does the stored coefficient match the number printed in the original paper? The agent is deliberately isolated from methodology. If the paper's specification is itself problematic, that is the methodology reviewers' concern, not the fidelity agent's.

The paper-auditor opens \texttt{pdf/[id].pdf}, locates the main-preferred-baseline specification using a four-level hierarchy (main $>$ preferred $>$ baseline $>$ column 1 of the headline table), extracts the coefficient, standard error, $N$, cluster level, and significance stars, and emits one of five verdicts.

\begin{table}[H]
\centering
\footnotesize
\begin{tabularx}{\textwidth}{@{}l X@{}}
\toprule
\textbf{Verdict} & \textbf{Criterion} \\
\midrule
EXACT             & Relative gap below 1\%. \\
WITHIN\_TOLERANCE & Relative gap between 1\% and 5\%. \\
WARN              & Relative gap between 5\% and 20\%, with documented cause. \\
FAIL              & Relative gap above 20\% or sign flip. \\
NOT\_APPLICABLE   & Paper does not report a comparable scalar (event-study-only or structurally different estimand). \\
\bottomrule
\end{tabularx}
\end{table}

Verdicts are written to \texttt{paper\_audit.md} per article and aggregated in \texttt{analysis/paper\_fidelity.csv}. The distribution over the 53-article sample is 36 EXACT, 5 WITHIN\_TOLERANCE, 4 WARN, and 8 NOT\_APPLICABLE.


\\subsection{Orchestration and credibility rating}
\label{sec:appc_skeptic}

The \texttt{skeptic} agent composes the audit layer into a single per-paper credibility report. It invokes the five methodology reviewers and the paper-auditor in parallel for one article and synthesises their verdicts into one overall rating on a three-axis rubric.

Axis~1 is Fidelity. It summarises the paper-auditor verdict. F-HIGH covers EXACT or WITHIN\_TOLERANCE. F-MOD covers WARN. F-LOW covers FAIL. F-NA covers NOT\_APPLICABLE.

Axis~2 is Implementation. It aggregates the five methodology reviewers. I-HIGH requires all PASS. I-MOD requires one WARN. I-LOW is assigned when there are two or more WARN or any FAIL.

Axis~3 is Design credibility. It reports D-ROBUST, D-MODERATE, D-FRAGILE, D-BROKEN, or D-NA as a separate finding about the paper's identification strategy, derived from HonestDiD $\bar M$ breakdown values, the Bacon TvT share, and pre-trend patterns. The design axis is a finding about the paper, not a demerit against the reanalysis.

The overall rating is the fixed product of Axes~1 and~2.

\begin{table}[H]
\centering
\footnotesize
\begin{tabularx}{\textwidth}{@{}l X@{}}
\toprule
\textbf{Rating} & \textbf{Axis combination} \\
\midrule
HIGH     & F-HIGH $\times$ I-HIGH, or F-NA $\times$ I-HIGH when no published scalar exists. \\
MODERATE & One axis degraded (one WARN on either axis). \\
LOW      & Multiple degradations or any implementation FAIL. \\
VERY LOW & Catastrophic failures (none in the final sample). \\
\bottomrule
\end{tabularx}
\end{table}

Decoupling Axis~3 from the rating is deliberate. A paper may be faithfully reproduced and cleanly implemented (HIGH) while still exhibiting fragile identification under modern diagnostics (D-FRAGILE). Under this rubric, design fragility does not reduce the credibility assigned to the reanalysis. It is reported as a distinct finding about the paper itself. This is the correct separation. The reanalysis is sound, and the paper has an identification concern.


\\subsection{Knowledge base and pattern-curator}
\label{sec:appc_knowledge}

Large-scale reanalysis is iterative, as \textcite{xu2026scaling} emphasize, because the system encounters novel implementation failures that must be diagnosed, resolved, and remembered. This project addresses that challenge through the version-controlled knowledge base \texttt{knowledge/failure\_patterns.md}, maintained by the \texttt{pattern-curator} agent. Each entry records the estimation step and data configuration that triggered the failure, the underlying technical reason, and a reusable corrective instruction.

The catalog contains 51 numbered entries, grouped in five families.

\begin{table}[H]
\centering
\footnotesize
\begin{tabularx}{\textwidth}{@{}l l X@{}}
\toprule
\textbf{Family} & \textbf{Patterns} & \textbf{Representative issues} \\
\midrule
Estimator output extraction            & 1--4, 19           & Sun--Abraham \texttt{iplot} missing SEs, Gardner term naming, negative horizons silently dropped, BJS incompatibility with individual-level RCS. \\
CS-DID estimation                      & 5, 8, 15, 22--26, 42--43, 49--51 & Singular covariances in small cohorts, RCS panel configuration, cohort variables that vary within units, base-period and cumulative-vs-marginal effects, late-treated absorption, propensity-score overfit (42), TWFE/CS sample mismatch (49--50), RCS control collinearity silently inflating ATT (51). \\
Stata-to-R data translation            & 7, 9--12, 16       & Half-year and monthly date formats, non-contiguous time periods, filter syntax differences, treatment-variable construction, \texttt{xi}-expanded variable names. \\
Event-study construction               & 13--14, 18, 20--21 & Tail-bin renaming inflating post-treatment average, pre-treatment cohort identification, never-treated unit exclusion, event-dummy collinearity with controls, sample restrictions. \\
TWFE specification and platform issues & 6, 17              & Plain OLS without unit and time FEs (systematic ${\sim}10\%$ divergence when the template imposes them), Windows-specific PDF file-locking. \\
\bottomrule
\end{tabularx}
\end{table}

The orchestrator reads this knowledge base before profiling each new article, so known pitfalls (Stata monthly dates, propensity-score overfit with many geographic dummies, and so on) are handled preemptively in \texttt{metadata.json} rather than discovered at the Analyst stage.


\\subsection{Repository custodian}
\label{sec:appc_custodian}

The \texttt{repo-custodian} agent is the single git and GitHub steward. It owns three things. The first is \texttt{code/audit/full\_sync.sh}, which refreshes the \texttt{overleaf/} and \texttt{health\_did\_replication/} LaTeX projects from the canonical \texttt{Replication Package/}. The second is \texttt{code/audit/run\_audit.R}, a 24-check verification battery that must pass before every commit. The third is all commits, pushes, and the verification that no raw data, secrets, or oversize files are staged.

All reviewer agents are strictly read-only. Only \texttt{repo-custodian} and \texttt{pattern-curator} may modify files under the canonical package, and both have narrowly scoped write permissions documented in their agent definitions.


\\subsection{The metadata contract}
\label{sec:appc_metadata}

The \texttt{metadata.json} file is the sole interface between the Profiler (scientific reasoning) and the Analyst (deterministic computation). Table~\ref{tab:metadata_fields} documents its key fields.

\begin{longtable}{>{\ttfamily\footnotesize\raggedright}p{6.0cm} >{\small\raggedright\arraybackslash}p{9.0cm}}
\caption{Key fields of \texttt{metadata.json}.} \label{tab:metadata_fields} \\
\toprule
\textnormal{\textbf{Field}} & \textbf{Description} \\
\midrule
\endfirsthead
\toprule
\textnormal{\textbf{Field}} & \textbf{Description} \\
\midrule
\endhead
\midrule
\multicolumn{2}{r}{\footnotesize\itshape Continued on next page} \\
\endfoot
\bottomrule
\endlastfoot

id & Mapping ID for the article \\
title & Full article title \\
group\_label & Category. One of \texttt{escalonada\_controles}, \texttt{escalonada\_sem\_controles}, \texttt{unica\_controles}, or \texttt{unica\_sem\_controles} \\
data.file & Path to the data file relative to project root \\
data.sample\_filter & R expression applied via \texttt{dplyr::filter()} \\[3pt]
variables.outcome & Outcome variable name \\
variables.unit\_id & Panel unit identifier \\
variables.time & Time-period variable \\
variables.treatment\_twfe & Binary post-treatment indicator for TWFE \\
variables.gvar\_cs & Cohort variable for CS-DID (0 = never-treated) \\
variables.twfe\_controls & Time-varying controls in the TWFE regression \\
variables.cs\_controls & Controls passed to \texttt{att\_gt(xformla=)} \\
variables.cluster & Clustering variable for standard errors \\
variables.weight & Observation weights \\
variables.additional\_fes & Extra fixed effects beyond unit + time \\[3pt]
panel\_setup.data\_structure & \texttt{"panel"} or \texttt{"repeated\_cross\_section"} \\
panel\_setup.treatment\_timing & \texttt{"staggered"} or \texttt{"single"} \\[3pt]
analysis.run\_csdid\_nyt & Run CS-DID with not-yet-treated controls \\
analysis.run\_bacon & Run Bacon decomposition \\
analysis.has\_event\_study & Whether the original paper reports an event study \\
analysis.event\_pre & Number of pre-treatment periods in the event window \\
analysis.event\_post & Number of post-treatment periods in the event window \\[3pt]
preprocessing.time\_recode & Mapping for non-contiguous time periods \\
preprocessing.construct\_vars & R expressions evaluated before estimation \\[3pt]
original\_result.beta\_twfe & Original TWFE coefficient (for validation) \\

\end{longtable}

The \texttt{preprocessing.construct\_vars} field deserves particular attention. Because the R template is never modified, any article-specific data manipulation (constructing a treatment indicator, creating a cohort variable, recoding missing values) must be expressed as R code strings in this array. These expressions are evaluated sequentially within the template's data environment before any estimation begins. A typical entry might define a post-treatment indicator as \verb|"df$Post_avg <- as.integer(df$year >= df$treat_year & df$treated == 1)"| or filter missing observations as \verb|"df <- df %>% filter(!is.na(df$outcome_var))"|. This mechanism keeps the template fixed while accommodating the heterogeneity of 53 different replication packages.


\\subsection{Calibration and validation}
\label{sec:appc_calibration}

The pipeline is validated against a calibration set of 10 articles for which the correct TWFE and CS-DID estimates were known in advance from prior manual reanalysis by the author. Two validation thresholds govern acceptance on any article. The reestimated TWFE coefficient must match the original paper's coefficient up to rounding. The reestimated CS-DID estimate must fall within 10\% of the manually obtained value, with the threshold allowing for minor R-vs-Stata implementation differences in \texttt{att\_gt}.

When an article fails either threshold, the failure is diagnosed, the \texttt{pattern-curator} appends a new entry to the knowledge base, and the orchestrator reruns the article with the updated rules. All 10 calibration articles pass both thresholds under the knowledge base in its current state. The remaining 43 articles are processed through the same pipeline. The knowledge base ensures that known pitfalls are handled preemptively in \texttt{metadata.json} rather than discovered at the Analyst stage.


\\subsection{Comparison with \texorpdfstring{\textcite{xu2026scaling}}{Xu and Yang (2026)}}
\label{sec:appc_xuyang}

Both frameworks separate scientific reasoning (what to estimate) from computational execution (how to estimate it), treating that division as the central architectural decision for scalable reanalysis. The main differences are summarized in Table~\ref{tab:framework_comparison}.

\begin{table}[H]
\centering
\caption{Framework comparison: \textcite{xu2026scaling} versus the present project.}
\label{tab:framework_comparison}
\small
\begin{tabular}{>{\raggedright\arraybackslash}p{3.4cm}
                >{\raggedright\arraybackslash}p{5.4cm}
                >{\raggedright\arraybackslash}p{5.4cm}}
\toprule
\textbf{Dimension} & \textbf{\textcite{xu2026scaling}} & \textbf{Present project} \\
\midrule
Scope                 & 92 IV studies (215 specifications)                                                   & 53 DiD studies \\[4pt]
Estimators            & 2SLS + first-stage F-stat                                                            & TWFE, CS-DID (NT/NYT), Sun--Abraham, Gardner, Bacon, HonestDiD \\[4pt]
Agents                & 7 (Profiler, Librarian, Janitor, Runner, Skeptic, Journalist, QC)                    & 9 (3 execution + 5 methodology reviewers + paper-auditor + skeptic + pattern-curator + repo-custodian) \\[4pt]
Interface contract    & Structured extraction fields                                                         & \texttt{metadata.json} (30+ fields) \\[4pt]
Execution layer       & Custom Stata/R scripts                                                               & Single fixed R template \\[4pt]
Fidelity verification & Final QC stage                                                                       & \texttt{paper-auditor} agent, 5-level verdict rubric \\[4pt]
Methodology audit     & Not separate                                                                         & Five dedicated read-only reviewers (TWFE, CS-DID, Bacon, HonestDiD, dChH) \\[4pt]
Credibility rating    & Not explicit                                                                         & Three-axis (Fidelity $\times$ Implementation gives HIGH/MOD/LOW, Design credibility separate finding) \\[4pt]
Knowledge base        & Not explicitly described                                                             & \texttt{failure\_patterns.md} (51 numbered entries) \\[4pt]
Calibration           & 67 manually verified studies                                                         & 10 articles with prior manual reanalysis \\[4pt]
Primary challenge     & Volume (92 articles)                                                                 & Heterogeneity across 53 packages, fidelity verification, and design-vs-implementation disentangling \\
\bottomrule
\end{tabular}
\end{table}

The primary distinction is scale versus depth. \textcite{xu2026scaling} process 92 articles with relatively homogeneous estimation (2SLS), whereas this project processes fewer articles with substantial heterogeneity in data formats, variable encoding, and panel structure. The fixed-template design with \texttt{construct\_vars} preprocessing is the response to that heterogeneity. All idiosyncratic data manipulation is expressed declaratively in \texttt{metadata.json}, keeping estimation code uniform across articles. The nine-agent architecture then adds a methodology audit layer, a dedicated fidelity layer, and an explicit credibility rubric not present in the original framework.


\\subsection{Directory structure}
\label{sec:appc_directory}

The canonical project directory is organized as follows.

\begin{tcolorbox}[colback=gray!5, colframe=gray!60, width=\textwidth]
\footnotesize
\begin{verbatim}
Replication Package/
  CLAUDE.md                                 # Master instructions (orchestrator)
  code/
    analysis/
      did_analysis_template.R               # Fixed R template
      01_run_all_did.R                      # Pipeline driver
    aggregation/
      01_consolidate_results.R              # Master results table
    tables/                                 # Table-generating scripts
    figures/                                # Figure-generating scripts
    audit/
      run_audit.R                           # 24-check verification
      full_sync.sh                          # Sync deposit -> LaTeX projects
  data/
    metadata/[id].json                      # Per-article specification
    raw/                                    # Replication data (junctioned)
  analysis/
    consolidated_results.csv                # Master table (all estimators)
    skeptic_ratings.csv                     # Per-paper 3-axis ratings
    paper_fidelity.csv                      # Per-paper fidelity verdicts
  knowledge/
    failure_patterns.md                     # 51 version-controlled patterns
  results/
    by_article/[id]/
      metadata.json
      results.csv
      event_study.pdf
      skeptic_report.md
      paper_audit.md
  output/
    figures/                                # Final figures (.pdf/.png)
    tables/                                 # Auto-generated .tex tables
  overleaf/                                 # Dissertation LaTeX source
  health_did_replication/                   # Working-paper LaTeX source
\end{verbatim}
\end{tcolorbox}


\\subsection{Replication instructions}
\label{sec:appc_replication}

Replicating or extending this project requires R (version $\geq 4.4$) with the packages \texttt{fixest}, \texttt{did}, \texttt{did2s}, \texttt{bacondecomp}, \texttt{HonestDiD}, \texttt{haven}, \texttt{dplyr}, \texttt{ggplot2}, \texttt{tidyr}, and \texttt{jsonlite}. It also requires access to a language model capable of reading PDFs and structured files. The orchestrator in this project is Claude Sonnet~4 (Anthropic, 2025), accessed through Claude Code, with the instruction set in \texttt{CLAUDE.md}.

To add a new article, place the article PDF in \texttt{pdf/[id].pdf} and the replication package in \texttt{data/raw/[id]/}, then create \texttt{results/by\_article/[id]/}. The Profiler, following \texttt{skills/SKILL\_profiler.md}, reads the article and replication code to fill \texttt{data/metadata/[id].json}. This step requires scientific judgment in identifying the main specification, mapping variable names, and determining whether adoption is staggered. The knowledge base (\texttt{failure\_patterns.md}) should be consulted for known pitfalls. The Analyst is invoked with
\begin{quote}
\texttt{Rscript code/analysis/did\_analysis\_template.R [id]}.
\end{quote}
If the script fails, the error is diagnosed, \texttt{metadata.json} corrected, and the script rerun. New failure modes are appended to the knowledge base by the \texttt{pattern-curator} using the standard Context, Root-Cause, Resolution-Rule schema. The Reporter then reads \texttt{results.csv} and \texttt{event\_study.pdf} to write \texttt{report.md}. Finally, the six audit agents are composed by the \texttt{skeptic} to produce the credibility rating in \texttt{skeptic\_report.md}.

For articles in the calibration set described in Section~\ref{sec:appc_calibration}, validation involves comparing the template output against the manual estimates. TWFE must match the original paper up to rounding. CS-DID may differ by up to 10\% due to R-vs-Stata implementation differences. Articles without prior manual analysis are validated through the audit layer. The \texttt{paper-auditor} verdict checks numerical fidelity, the five methodology reviewers check implementation, and the \texttt{skeptic} synthesises the result.
